\section{数值计算}

\subsection{多项式与线性递推}
  \kactlimport{Polynomial.h}
  % \kactlimport{PolyRoots.h}
  \kactlimport{PolyInterpolate.h}
  \kactlimport{LagrangeInterp.h}
  \kactlimport{BerlekampMassey.h}
  \kactlimport{LinearRecurrence.h}

\subsection{最优化}
  \kactlimport{GoldenSectionSearch.h}
  % \kactlimport{HillClimbing.h}
  % \kactlimport{Integrate.h}
  % \kactlimport{IntegrateAdaptive.h}
  % \kactlimport{Simplex.h}

\subsection{矩阵}
  \kactlimport{Determinant.h}
  \kactlimport{IntDeterminant.h}
  \kactlimport{SolveLinear.h}
  % \kactlimport{SolveLinear2.h}
  \kactlimport{SolveLinearBinary.h}
  \kactlimport{MatrixInverse.h}
  % \kactlimport{MatrixInverse-mod.h}
  % \kactlimport{Tridiagonal.h}

\subsection{FFT, NTT, FWT}
	\kactlimport{FastFourierTransform.h}
	\kactlimport{FastFourierTransformMod.h}
	\kactlimport{NumberTheoreticTransform.h}
	\kactlimport{FWT.h}
  \kactlimport{MaxPlusConvolution.h}

\paragraph{卷积用法}
下面的代码直接调用上面的 \texttt{conv} (FFT), \texttt{nttConv} (NTT), \texttt{xorConv} 等 (FWT). 下标都从 0 开始.
\begin{compactitem}
\item \textbf{多项式乘法 / 两数组卷积}: $c_k = \sum_i a_i b_{k-i}$, 长度 $|a|+|b|-1$. 差卷积 $\sum_i a_i b_{i+k}$ 把 $b$ 反转.
\begin{lstlisting}
// a[i], b[i] 为 x^i 的系数
vector<ll> c=nttConv(a,b); // 模 998244353, 先取模
vector<ll> d=conv(a,b); // 精确整数 (FFT), 可有负数
// 差卷积 e[k]=sum a[i]*b[i+k], k 从 -(n-1) 到 m-1:
vector<ll> rb(b.rbegin(),b.rend()); // m=b.size()
vector<ll> t=conv(a,rb); // e[k] 即 t[m-1-k]
\end{lstlisting}
\item \textbf{大整数乘法}: 低位在前存数字, 卷积后进位.
\begin{lstlisting}
string bigMul(string s,string t){ // 十进制, 无负号
	int n=s.size();
	int m=t.size();
	vector<ll> a(n),b(m); // 低位在前
	for(int i=0;i<n;i++)
		a[i]=s[n-1-i]-'0';
	for(int i=0;i<m;i++)
		b[i]=t[m-1-i]-'0';
	vector<ll> c=conv(a,b);
	string r;
	ll carry=0;
	for(int i=0;i<(int)c.size()||carry;i++){
		if(i<(int)c.size())
			carry+=c[i];
		r+=char('0'+carry%10);
		carry/=10;
	}
	while(r.size()>1&&r.back()=='0') // 去前导 0
		r.pop_back();
	reverse(r.begin(),r.end());
	return r;
}
\end{lstlisting}
\item \textbf{两数之和计数}: 求 $v_i + v_j = k$ ($i < j$) 的对数, 值域 $[0, V]$. 自己卷自己得到有序对且含 $i = j$.
\begin{lstlisting}
vector<ll> pairSum(vector<int>&v,int V){
	vector<ll> cnt(V+1);
	for(int x:v)
		cnt[x]++;
	vector<ll> f=conv(cnt,cnt); // 有序对, 含 i=j
	for(int x:v)
		f[2*x]--; // 去掉 i=j
	for(auto&y:f)
		y/=2; // 变成无序对 i<j
	return f; // f[k] 为对数, 长度 2V+1
}
\end{lstlisting}
\item \textbf{带通配符的字符串匹配}: \texttt{*} 匹配任意字符. 字符值 $1..26$, 通配符为 0,
位置 $i$ 匹配当且仅当 $\sum_j (p_j - s_{i+j})^2 p_j s_{i+j} = 0$. 展开后是三个卷积 ($p$ 反转).
NTT 下和可能恰为模数的倍数, 误判概率极小; 介意可再换一个模数验证.
\begin{lstlisting}
ll wv(char c){return c=='*'?0:c-'a'+1;} // 通配符为 0
vector<int> wildMatch(string s,string p){
	int n=s.size();
	int m=p.size();
	vector<int> res; // 匹配的起点, 下标从 0 开始
	if(m>n)
		return res;
	vector<ll> S[4],P[4]; // S[k][i] 为 s_i 的 k 次方
	for(int k=1;k<=3;k++){
		S[k].resize(n);
		P[k].resize(m);
	}
	for(int i=0;i<n;i++){
		ll v=wv(s[i]);
		S[1][i]=v;
		S[2][i]=v*v;
		S[3][i]=v*v*v;
	}
	for(int j=0;j<m;j++){ // P 为 p 反转
		ll v=wv(p[m-1-j]);
		P[1][j]=v;
		P[2][j]=v*v;
		P[3][j]=v*v*v;
	}
	vector<ll> c1=nttConv(P[3],S[1]);
	vector<ll> c2=nttConv(P[2],S[2]);
	vector<ll> c3=nttConv(P[1],S[3]);
	for(int i=0;i+m<=n;i++){
		int t=i+m-1;
		if((c1[t]-2*c2[t]+c3[t]+2*MOD)%MOD==0)
			res.push_back(i);
	}
	return res;
}
\end{lstlisting}
\item \textbf{位运算卷积}: 求 $v_i \oplus v_j = k$ ($i < j$) 的对数 (值在 $[0, 2^K)$), 以及子集和.
\begin{lstlisting}
vector<ll> pairXor(vector<int>&v,int K){
	vector<ll> cnt(1<<K);
	for(int x:v)
		cnt[x]++;
	vector<ll> f=xorConv(cnt,cnt); // 有序对, 含 i=j
	f[0]-=v.size(); // i=j 时 xor 为 0
	for(auto&y:f)
		y/=2;
	return f;
}
// or 同理: f=orConv(cnt,cnt), 对每个 x 做 f[x]--, 再除 2
// 子集和: g[S]=sum a[T], T 取遍 S 的子集
vector<ll> g=a;
fwtOr(g,1); // fwtOr(g,-1) 还原
// 超集和: fwtAnd(g,1)
\end{lstlisting}
\end{compactitem}
